\documentclass[10pt,a4paper]{amsart}
\usepackage{amsmath,amssymb,amsthm,graphicx}
% For commutative diagrams we use Xy-pic
\usepackage[all]{xy}

% We use 2cell for 2-commutative diagrams.
\xyoption{2cell}
% Theorem environments.
%
\theoremstyle{plain}
\newtheorem{theorem}[subsection]{Theorem}
\newtheorem{proposition}[subsection]{Proposition}
\newtheorem{lemma}[subsection]{Lemma}

\theoremstyle{definition}
\newtheorem{definition}[subsection]{Definition}
\newtheorem{example}[subsection]{Example}
\newtheorem{exercise}[subsection]{Exercise}
\newtheorem{situation}[subsection]{Situation}

\theoremstyle{remark}
\newtheorem{remark}[subsection]{Remark}
\newtheorem{remarks}[subsection]{Remarks}

\numberwithin{equation}{subsection}

% Macros
%
\def\lim{\mathop{\mathrm{lim}}\nolimits}
\def\colim{\mathop{\mathrm{colim}}\nolimits}
\def\Spec{\mathop{\mathrm{Spec}}}
\def\Hom{\mathop{\mathrm{Hom}}\nolimits}
\def\Ext{\mathop{\mathrm{Ext}}\nolimits}
\def\SheafHom{\mathop{\mathcal{H}\!\mathit{om}}\nolimits}
\def\SheafExt{\mathop{\mathcal{E}\!\mathit{xt}}\nolimits}
\def\Sch{\mathit{Sch}}
\def\Mor{\mathop{\mathrm{Mor}}\nolimits}
\def\Ob{\mathop{\mathrm{Ob}}\nolimits}
\def\Sh{\mathop{\mathit{Sh}}\nolimits}
\def\NL{\mathop{N\!L}\nolimits}
\def\CH{\mathop{\mathrm{CH}}\nolimits}
\def\proetale{{pro\text{-}\acute{e}tale}}
\def\etale{{\acute{e}tale}}
\def\QCoh{\mathit{QCoh}}
\def\Ker{\mathop{\mathrm{Ker}}}
\def\Im{\mathop{\mathrm{Im}}}
\def\Coker{\mathop{\mathrm{Coker}}}
\def\Coim{\mathop{\mathrm{Coim}}}

% Boxtimes
%
\DeclareMathSymbol{\boxtimes}{\mathbin}{AMSa}{"02}

%
% Macros for moduli stacks/spaces
%
\def\QCohstack{\mathcal{QC}\!\mathit{oh}}
\def\Cohstack{\mathcal{C}\!\mathit{oh}}
\def\Spacesstack{\mathcal{S}\!\mathit{paces}}
\def\Quotfunctor{\mathrm{Quot}}
\def\Hilbfunctor{\mathrm{Hilb}}
\def\Curvesstack{\mathcal{C}\!\mathit{urves}}
\def\Polarizedstack{\mathcal{P}\!\mathit{olarized}}
\def\Complexesstack{\mathcal{C}\!\mathit{omplexes}}
% \Pic is the operator that assigns to X its picard group, usage \Pic(X)
% \Picardstack_{X/B} denotes the Picard stack of X over B
% \Picardfunctor_{X/B} denotes the Picard functor of X over B
\def\Pic{\mathop{\mathrm{Pic}}\nolimits}
\def\Picardstack{\mathcal{P}\!\mathit{ic}}
\def\Picardfunctor{\mathrm{Pic}}
\def\Deformationcategory{\mathcal{D}\!\mathit{ef}}
\usepackage{cleveref}
\crefname{originallemma}{original lemma}{original lemmas}
\Crefname{originallemma}{Original lemma}{Original lemmas}
\makeatletter
\newlabel{morphisms-lemma-etale-locally-quasi-finite}{{03WS}{source}{}{sourcepointer-1}{}}
\newlabel{morphisms-lemma-etale-locally-quasi-finite@cref}{{[originallemma][0][]03WS}{[1][0][]source}{}{}{}}
\makeatother
\crefname{originallemma}{original lemma}{original lemmas}
\Crefname{originallemma}{Original lemma}{Original lemmas}
\makeatletter
\newlabel{morphisms-lemma-composition-quasi-finite}{{01TL}{source}{}{sourcepointer-2}{}}
\newlabel{morphisms-lemma-composition-quasi-finite@cref}{{[originallemma][0][]01TL}{[1][0][]source}{}{}{}}
\makeatother
\crefname{originallemma}{original lemma}{original lemmas}
\Crefname{originallemma}{Original lemma}{Original lemmas}
\makeatletter
\newlabel{more-morphisms-lemma-etale-splits-off-quasi-finite-part-technical-variant}{{02LO}{source}{}{sourcepointer-3}{}}
\newlabel{more-morphisms-lemma-etale-splits-off-quasi-finite-part-technical-variant@cref}{{[originallemma][0][]02LO}{[1][0][]source}{}{}{}}
\makeatother
\crefname{originallemma}{original lemma}{original lemmas}
\Crefname{originallemma}{Original lemma}{Original lemmas}
\makeatletter
\newlabel{etale-lemma-finite-etale-one-point}{{04DH}{source}{}{sourcepointer-4}{}}
\newlabel{etale-lemma-finite-etale-one-point@cref}{{[originallemma][0][]04DH}{[1][0][]source}{}{}{}}
\makeatother
\title{A separated locally quasi-finite morphism has property A}
\author{The Stacks Project Authors}
\date{stacks-project:etale-cohomology.tex:local-snapshot-20261011; GNU FDL 1.2 or later}
\begin{document}
\maketitle
\noindent\textit{Editorial note.} All delivered mathematical text is the author's own, copied from the pinned Stacks Project chapter file etale-cohomology.tex; no statement or proof step was rewritten, reordered or removed.\par\medskip
\noindent\textit{Editorial note.} The proof is the author's complete argument for the property, and the construction of the etale neighbourhood and the decomposition is carried out inside the counted span.\par\medskip
\noindent\textit{Editorial note.} Every labelled reference occurring inside the counted spans is carried in assembly.external\_references as a non-load-bearing navigation pointer (same-chapter ones at their definition site, cross-chapter ones at the invocation inside the counted span), so no load-bearing mathematics is replaced by a stub. No bibliography citation occurs inside the counted spans.\par\medskip
\noindent\textit{Editorial note.} The checked chapter list for same-name collisions is the set of stacks chapters staged in this repository: chow, cohomology, divisors, etale-cohomology, flat, more-algebra, more-morphisms, morphisms, spaces-more-morphisms, varieties together with this card's chapter etale-cohomology.\par\medskip
\noindent\textit{Editorial note.} The preamble is the chapter's own native preamble excerpt already used for this source: the Xy-pic setup (native preamble lines 14-18) and the native theorem declarations with the mathematical macros (native preamble lines 154-217).\par\medskip
\noindent\textit{Allowed standard prerequisites (editorial).}\begin{itemize}
\item Property (A) for a morphism of schemes. Same chapter, the formulation of property (A) delivered verbatim as a context span in this card (source lines 5708-5726). Supplies the condition verified for a separated locally quasi-finite morphism.
\item Separation and local quasi-finiteness of a morphism, and the openness of the image of an etale morphism. Same chapter together with the chapter Morphisms, the consequences of separation and local quasi-finiteness used to isolate a point in its fibre and to spread it out over the base. Supplies the construction of the etale neighbourhood V of Y and of the disjoint union decomposition.
\end{itemize}\medskip
\ifdefined\hypertarget\hypertarget{sourcepointer-1}{}\fi
\noindent\textit{Source reference.} Original lemma 03WS: sources/stacks-etale-cohomology/pinned-original-source.tex, source lines 5806--5806. Navigation pointer invoked at proof line 5806; the counted span contains the pointer only and the cited result is not reproduced inside it, so the pointer is registered as a non-load-bearing navigation stub in the established form for Stacks material. No bibliography citation is involved.\par\medskip
\ifdefined\hypertarget\hypertarget{sourcepointer-2}{}\fi
\noindent\textit{Source reference.} Original lemma 01TL: sources/stacks-etale-cohomology/pinned-original-source.tex, source lines 5808--5808. Navigation pointer invoked at proof line 5808; the counted span contains the pointer only and the cited result is not reproduced inside it, so the pointer is registered as a non-load-bearing navigation stub in the established form for Stacks material. No bibliography citation is involved.\par\medskip
\ifdefined\hypertarget\hypertarget{sourcepointer-3}{}\fi
\noindent\textit{Source reference.} Original lemma 02LO: sources/stacks-etale-cohomology/pinned-original-source.tex, source lines 5811--5811. Navigation pointer invoked at proof line 5811; the counted span contains the pointer only and the cited result is not reproduced inside it, so the pointer is registered as a non-load-bearing navigation stub in the established form for Stacks material. No bibliography citation is involved.\par\medskip
\ifdefined\hypertarget\hypertarget{sourcepointer-4}{}\fi
\noindent\textit{Source reference.} Original lemma 04DH: sources/stacks-etale-cohomology/pinned-original-source.tex, source lines 5843--5843. Navigation pointer invoked at proof line 5843; the counted span contains the pointer only and the cited result is not reproduced inside it, so the pointer is registered as a non-load-bearing navigation stub in the established form for Stacks material. No bibliography citation is involved.\par\medskip
\noindent
Let $f : X \to Y$ be a morphism of schemes.
Here is a list of conditions we will consider in the following:
\begin{enumerate}
\item[(A)] For every \'etale morphism $U \to X$ and $u \in U$ there exist
an \'etale morphism $V \to Y$ and a disjoint union decomposition
$X \times_Y V = W \amalg W'$ and a morphism $h : W \to U$ over $X$
with $u$ in the image of $h$.
\item[(B)] For every $V \to Y$ \'etale, and every \'etale covering
$\{U_i \to X \times_Y V\}$ there exists an \'etale covering
$\{V_j \to V\}$ such that for each $j$ we have
$X \times_Y V_j = \coprod W_{ij}$ where $W_{ij} \to X \times_Y V$
factors through $U_i \to X \times_Y V$ for some $i$.
\item[(C)] For every $U \to X$ \'etale, there exists a $V \to Y$ \'etale
and a surjective morphism $X \times_Y V \to U$ over $X$.
\end{enumerate}
It turns out that each of these properties has meaning in terms of
the behaviour of the functor $f_{small, *}$. We will work this
out in the next few sections.

\begin{lemma}
\label{lemma-locally-quasi-finite-A}
Let $f : X \to Y$ be a separated locally quasi-finite morphism of schemes.
Then property (A) above holds.
\end{lemma}

\begin{proof}
Let $U \to X$ be an \'etale morphism and $u \in U$.
The geometric statement (A) reduces directly to the case where $U$ and $Y$
are affine schemes. Denote $x \in X$ and $y \in Y$ the
images of $u$. Since $X \to Y$ is locally quasi-finite, and $U \to X$ is
locally quasi-finite (see
Morphisms, Lemma \ref{morphisms-lemma-etale-locally-quasi-finite})
we see that $U \to Y$ is locally quasi-finite (see
Morphisms, Lemma \ref{morphisms-lemma-composition-quasi-finite}).
Moreover both $X \to Y$ and $U \to Y$ are separated. Thus
More on Morphisms, Lemma
\ref{more-morphisms-lemma-etale-splits-off-quasi-finite-part-technical-variant}
applies to both morphisms. This means we may pick an \'etale neighbourhood
$(V, v) \to (Y, y)$ such that
$$
X \times_Y V = W \amalg R, \quad
U \times_Y V = W' \amalg R'
$$
and points $w \in W$, $w' \in W'$ such that
\begin{enumerate}
\item $W$, $R$ are open and closed in $X \times_Y V$,
\item $W'$, $R'$ are open and closed in $U \times_Y V$,
\item $W \to V$ and $W' \to V$ are finite,
\item $w$, $w'$ map to $v$,
\item $\kappa(v) \subset \kappa(w)$ and $\kappa(v) \subset \kappa(w')$
are purely inseparable, and
\item no other point of $W$ or $W'$ maps to $v$.
\end{enumerate}
Here is a commutative diagram
$$
\xymatrix{
U \ar[d] & U \times_Y V \ar[l] \ar[d] & W' \amalg R' \ar[d] \ar[l] \\
X \ar[d] & X \times_Y V \ar[l] \ar[d] & W \amalg R \ar[l] \\
Y & V \ar[l]
}
$$
After shrinking $V$ we may assume that $W'$ maps into $W$:
just remove the image the inverse image of $R$ in $W'$; this is
a closed set (as $W' \to V$ is finite) not containing $v$.
Then $W' \to W$ is finite because both $W \to V$ and $W' \to V$ are finite.
Hence $W' \to W$ is finite \'etale, and there is exactly one point in the
fibre over $w$ with $\kappa(w) = \kappa(w')$. Hence $W' \to W$ is an
isomorphism in an open neighbourhood $W^\circ$ of $w$, see
\'Etale Morphisms, Lemma \ref{etale-lemma-finite-etale-one-point}.
Since $W \to V$ is finite the image of $W \setminus W^\circ$ is a closed
subset $T$ of $V$ not containing $v$. Thus after replacing $V$ by
$V \setminus T$ we may assume that $W' \to W$ is an isomorphism.
Now the decomposition $X \times_Y V = W \amalg R$ and the morphism
$W \to U$ are as desired and we win.
\end{proof}

\end{document}
